\subsection{Fourier transform of a promeasure}

Let E be a locally convex space and \(\mu = (\mu_{\mathrm{V}})_{\mathrm{V} \in \mathcal{F}(\mathrm{E})}\) a promeasure on E. For every continuous linear form \(x'\) on E, we denote by \(\mu_{x'}\) the measure on R that is the image under \(x'\) of the promeasure \(\mu\) on E. The Fourier transform of \(\mu\) is the function \(F\mu\) on \(E'\) defined by

\[
(\mathcal {F} \mu) (x ^ {\prime}) = \int_ {\mathbf {R}} e ^ {i t} d \mu_ {x ^ {\prime}} (t).\tag{2}
\]

Let \(\lambda\) be a bounded measure on E. The Fourier transform of \(\lambda\) is the function on \(\mathbf{E}'\) defined by

\[
(\mathcal {F} \lambda) (x ^ {\prime}) = \int_ {\mathrm{E}} e ^ {i \langle x, x ^ {\prime} \rangle} d \lambda (x).\tag{3}
\]

Let \(\mu\) be the promeasure associated with \(\lambda\) . For every \(x' \in E'\) , the measure \(\mu_{x'}\) on \(\mathbf{R}\) is the image under \(x': \mathrm{E} \to \mathbf{R}\) of the measure \(\lambda\) on \(\mathrm{E}\) ; from the formulas (2) and (3), one immediately deduces \(\mathcal{F}\mu = \mathcal{F}\lambda\) .

Let \(\mu\) be any promeasure on \(\mathbf{E}\) , and \(u\) a continuous linear mapping of \(\mathbf{E}\) into a locally convex space \(\mathbf{E}_1\) . Denote by \(^t u\) the linear mapping of \(\mathbf{E}_1'\) into \(\mathbf{E}'\) that is the transpose of \(u\) , and by \(\nu\) the promeasure \(u(\mu)\) on \(\mathbf{E}_1\) . For every \(x_1' \in \mathbf{E}_1'\) , we have \(^t u(x_1') = x_1' \circ u\) , whence

\[
\nu_ {x _ {1} ^ {\prime}} = x _ {1} ^ {\prime} (\nu) = x _ {1} ^ {\prime} (u (\mu)) = \left(x _ {1} ^ {\prime} \circ u\right) (\mu) = \mu^ {t} u \left(x _ {1} ^ {\prime}\right).
\]

Consequently,

\[
\mathscr {F} (u (\mu)) = (\mathscr {F} \mu) \circ {} ^ {t} u.\tag{4}
\]

In particular, let us take for u the canonical mapping \(p_{V}\) of E onto E/V (for \(V \in \mathcal{F}(E)\) ). The promeasure \(p_{\mathrm{V}}(\mu)\) on E/V is associated with the measure \(\mu_{V}\) , and \(^{t}p_{V}\) is an isomorphism of the dual of E/V onto the subspace \(V^{\circ}\) of \(E'\) orthogonal to V. If \((\mathrm{E}/\mathrm{V})'\) is identified with \(V^{\circ}\) by means of \(^{t}p_{V}\) , then

\[
(\mathcal {F} \mu) (x ^ {\prime}) = \int_ {\mathrm{E} / \mathrm{V}} e ^ {i \langle x, x ^ {\prime} \rangle} d \mu_ {\mathrm{V}} (x)\tag{5}
\]

for all \(x' \in V^{\circ}\) . One has \(E' = \bigcup_{V \in \mathcal{F}(E)} V^{\circ}\) , so that the preceding formula characterizes the function \(\mathcal{F}\mu\) on \(E'\) . Finally, if one sets \(x' = 0\) in (5), one sees that the total mass of \(\mu\) is equal to \((\mathcal{F}\mu)(0)\) .

PROPOSITION 3. --- Let E be a locally convex space. The mapping \(\mu \mapsto \mathcal{F}\mu\) of the set of promeasures on E into the set of functions on E' is injective.

The formula (5) permits reducing to the case that E is finite-dimensional; since every finite-dimensional space is isomorphic to a space \(R^{n}\) , we can even suppose that there exists an integer \(n \geqslant 0\) such that \(E = R^{n}\) . We therefore have to prove that if \(\mu\) is a bounded measure (not necessarily positive) on \(R^{n}\) and if

\[
\int_ {\mathbf {R} ^ {n}} e ^ {i \langle x, y \rangle} d \mu (x) = 0
\]

for every linear form \(y\) on \(\mathbf{R}^n\) , then \(\mu = 0\) .

For every integer \(m \geqslant 0\) , let \(\mathbf{G}_m\) be the subgroup \(m \cdot \mathbf{Z}^n\) of \(\mathbf{R}^n\) . Denote by \(\mathcal{C}_m\) the vector space of continuous functions \(f\) on \(\mathbf{R}^n\) such that \(f(x + a) = f(x)\) for \(x \in \mathbf{R}^n\) and \(a \in \mathbf{G}_m\) . By Prop. 8 of GT, X, §4, No.~4, every function in \(\mathcal{C}_m\) is the uniform limit of finite linear combinations of functions of the type \(x \mapsto e^{2\pi i\langle x,q\rangle}\) with \(q \in m^{-1} \cdot \mathbf{Z}^n\) . Therefore \(\mu(f) = 0\) for every function \(f \in \mathcal{C}_m\) .

Let \(f\) be a continuous function on \(\mathbf{R}^n\) with compact support. For every integer \(m \geqslant 0\) , set \(f_m(x) = \sum_{q \in G_m} f(x + q)\) . It is immediate that for every \(x \in \mathbf{R}^n\) , the preceding series has only finitely many terms, and that \(f_m\) belongs to \(\mathcal{C}_m\) . Moreover, it is easy to see that the sequence \((f_m)\) tends to \(f\) uniformly on every compact set, and that there exists a constant \(C \geqslant 0\) such that \(|f_m| \leqslant C\) for all \(m\) . Consequently, \(\mu(f) = \lim_{m \to \infty} \mu(f_m)\) by Prop. 12 of §5, No.~6. Since \(f_m \in \mathcal{C}_m\) , we have \(\mu(f_m) = 0\) , whence finally \(\mu(f) = 0\) . Thus \(\mu = 0\) .

*Remark. --- When E is finite-dimensional, every character of E is of the form \(x \mapsto e^{i\langle x,x' \rangle}\) with \(x' \in \mathrm{E}'\) (Théor. spect., Ch. II, §1, No.~9, Cor. 3 of Prop. 12). In this case, Prop. 3 follows from the uniqueness theorem for the Fourier transformation (loc. cit., §1, No.~6, Cor. of Prop. 6).*

